\scanpage{89}

% L1–5
D) Sous les conditions de B), soit \deleted{\unsure{$G$ de}} $\underline g$ le groupe des autom de $(X,\deleted{(\unsure{\tau_*|\tau_i})})$, i.e. des $G_n$-autom de $X$, où $G_n$ est le groupe engendré par les gén. $\underline\sigma_i$ ($0\leq i\leq n$), liés par relations $\underline\sigma_i\underline\sigma_j=\underline\sigma_j\underline\sigma_i$ si $i,j$ non consécutifs. On a \deleted{donc}

% L6
\[
\underline g\longrightarrow\underline{\operatorname{Aut}}\bigl(D_0,\overset{\text{id sur}}{(c\sqcup(d))}\bigr)
=\operatorname{Aut}\bigl(D_0,\overset{\text{id sur}}{(\Delta_0)}\bigr)
\]

% L5–7, interline right
\margin{produit de div. relatifs}

% L6–8, right
\margin{où $\Delta_0$ est le diviseur sur $D_0$ défini par les sections $c,d$}

% L6–10, left
\margin{\uline{NB} $\underline g$ est lisse \unsure{comm.|connexe|commutatif} rel. \unsure{conn.|un.} de dim rel $1$, ses fibres sont isomorphes à $\mathbb G_m$ ($c\ne d$) ou $\mathbb G_a$ ($c=d$).}

% L7–12
Ceci dit, cet hom est un \uline{isom} -- par suite $D_0^{**}\inserted{=D_0\smallsetminus\Delta_0}$ est un $\underline g$-torseur : $\exists$ loc sections (c'est le contenu de B,C), et deux sections sont conjuguées par un unique \inserted{$G_n$-}autom. Donc \deleted{\unsure{en|En}} en vertu de C, itou pour le schém. des \inserted{$(\tau_*)$-}drapeaux non dégénérés, qui est un torseur sous $\underline g$.

% L13–19
\uline{Scholie} Un système $(\tau_*,X_*)$ \inserted{non dégénéré} peut être considéré comme une \uline{rigidification} d'un système $(\tau_*)$ non dég, i.e. d'une représentation projective n. dég. de $G_n$. La cat. (sur $S$ fixé) des $(\tau_*\deleted{\unsure{,X_*}})$ non dég. est équivalente à celle des $(\tau_*,X_*)$ (i.e. des \uline{réal. géom non dég de $G_n$}) munis d'un torseur sous le groupe \inserted{$\unsure{g_0|\underline g_0}$} $\underline{\operatorname{Aut}}(D_0,\overset{\text{id sur}}{\unsure{c_0;d_0|\Delta_0}})$ correspondant à cette dernière\ldots\ Utilisant la section $s_0$ disjointe de $d$, \deleted{comme \unsure{et} $d$ comme} \inserted{le normant pt $1$, $c$ comme $0$ et $d$ comme $\infty$}

% L20–24
(si $c\ne d$ en tte fibre i.e. le centre n'est à l'$\infty$ sur aucune fibre) on trouve que ce groupe est $G_m$, donc \deleted{il s'agit} \inserted{dans ce} \deleted{de} \inserted{cas} \deleted{\unsure{don|la}} la cat. des $(\tau_*)$ non dég \unsure{non|son|avec} centre partout à dist. finie équivaut à la cat produit de la cat des $(\tau_*,X_*)$ par celle des Mod. inv. sur $S$.

% L25–26
E) Soit $\Pi$ un quasi-polyèdre combinatoire régulier. $G$ le groupe de ses autom, $R$ l'ens. des repères, qui

% archive sheet number: 88
